%=====================================================================
% 光学（02）外教 Lecture 3 · Polarization Devices & Interference
% 精讲讲义（中文讲解 + LaTeX 原生公式）
% 来源：subjects/02-光学/课件/Polarization Devices 1.pdf（110 页）
% 编译：xelatex -halt-on-error 两遍
%=====================================================================
\documentclass[11pt]{ctexart}
%=====================================================================
% theme-violet.tex —— 讲义精讲 · 紫花主题
% 用法：\documentclass[11pt]{ctexart}
%       \input{assets/theme-violet}
%       \begin{document} ... \end{document}
% 只用 ctexart 自带能力 + xcolor/enumitem/fancyhdr/pifont/amsmath，
% 不依赖 tikz/tcolorbox，保证 MiKTeX 可直接编译。
%=====================================================================

\usepackage[a4paper,left=2.0cm,right=2.0cm,top=2.1cm,bottom=2.3cm,%
            headsep=12pt,footskip=22pt]{geometry}
\usepackage{amsmath,amssymb}
\usepackage{needspace}
\usepackage{xcolor}
\usepackage{array}
\usepackage{enumitem}
\usepackage{pifont}
\usepackage{fancyhdr}
\usepackage[hidelinks]{hyperref}

%---------------------------------------------------------------- 配色
\definecolor{violet}{RGB}{107,78,143}      % 主色：紫
\definecolor{violetdeep}{RGB}{63,44,92}    % 深紫：标题/强调
\definecolor{petal}{RGB}{183,159,208}      % 花瓣：装饰线、花饰
\definecolor{bandviolet}{RGB}{247,243,251} % 浅紫底：档案/提示框
\definecolor{bandcream}{RGB}{252,249,242}  % 米白底：易错/注意框
\definecolor{rulepurple}{RGB}{217,207,230} % 分隔线
\definecolor{inkgray}{RGB}{92,96,110}      % 次要文字

%---------------------------------------------------------------- 字体
% 有 Noto CJK 用 Noto，没有就退回 ctex 默认（Windows 下为宋体/黑体）
\IfFontExistsTF{Noto Serif CJK SC}{\setCJKmainfont{Noto Serif CJK SC}}{}
\IfFontExistsTF{Noto Sans CJK SC}{\setCJKsansfont{Noto Sans CJK SC}}{}

%---------------------------------------------------------------- 版式
\setlist{nosep,leftmargin=1.9em}
\linespread{1.28}
\setlength{\parskip}{3.5pt}
\setlength{\parindent}{0pt}
\setlength{\tabcolsep}{5pt}
\renewcommand{\arraystretch}{1.28}

%---------------------------------------------------------------- 花饰
% \ding{95} = 五瓣小花，\ding{96} = 星形花；都来自 pifont，必定存在
\newcommand{\orn}{\textcolor{petal}{\ding{95}}}          % 行内小标
\newcommand{\ornstar}{\textcolor{petal}{\ding{96}}}      % 星形花
\newcommand{\flowerrule}{%                                % 居中花饰分隔线
  \par\vspace{7pt}%
  \noindent\makebox[\linewidth]{\color{petal}%
    \rule{0.20\linewidth}{0.6pt}\hspace{8pt}{\small\ding{95}}\hspace{4pt}%
    \raisebox{1.1pt}{\tiny\ding{95}}\hspace{4pt}{\small\ding{95}}\hspace{8pt}%
    \rule{0.20\linewidth}{0.6pt}}%
  \par\vspace{7pt}}

%---------------------------------------------------------------- 标题与节头
\newcommand{\doctitle}[3]{%   #1 主标题  #2 副标题  #3 元信息（可用 \\ 分行）
  \begin{center}
    {\color{petal}\small \orn\enspace #2\enspace \ornstar}\\[7pt]
    {\LARGE\bfseries\color{violetdeep} #1}\\[10pt]
    {\color{petal}\rule{0.55\linewidth}{1pt}}\\[4pt]
    {\small\color{inkgray} #3}
  \end{center}\vspace{2pt}}

\newcommand{\hdr}[2]{%   #1 章节号（如 §0）  #2 标题
  \par\vspace{12pt}%
  \noindent\colorbox{violetdeep}{\color{white}\bfseries\small\;#1\;}\hspace{7pt}%
  {\color{violetdeep}\large\bfseries #2}\par\nobreak
  \vspace{3pt}{\color{petal}\hrule height 1.2pt}\vspace{6pt}}

\newcommand{\subhdr}[1]{%
  \par\vspace{6pt}%
  \noindent{\color{violet}\bfseries #1}\par\vspace{1pt}}

%---------------------------------------------------------------- 提示框
\newcommand{\callout}[1]{% 浅紫底：讲义档案、读法说明
  \par\medskip\noindent
  \fcolorbox{rulepurple}{bandviolet}{\parbox{\dimexpr\linewidth-2\fboxsep-2\fboxrule}{\vspace{2pt}#1\vspace{2pt}}}%
  \par\medskip}

\newcommand{\warnbox}[1]{% 米白底：易错、边界条件
  \par\medskip\noindent
  \fcolorbox{rulepurple}{bandcream}{\parbox{\dimexpr\linewidth-2\fboxsep-2\fboxrule}{\vspace{2pt}#1\vspace{2pt}}}%
  \par\medskip}

\newcommand{\gapitem}[2]{% #1 缺口类型（紫） #2 内容
  \par\vspace{2pt}\noindent
  \hangindent=1.4em\hangafter=1
  {\color{petal}\small\ding{96}}\enspace{\bfseries\color{violetdeep}#1}\enspace #2\par}

% 题干块（\qstem）：每道例题／作业的「详解」开头必须先复述题干与出处。
% 起因：2026-10-25 第五章讲义把七道作业的题面集中放在一页、详解另起一页，
% 读者翻到「第 4 题」时看不到条件，等于一道没有题干的题目。
% 规则：详解块内自带题干，读者不翻页就能读懂题在问什么。
\newcommand{\qstem}[1]{%
  \par\medskip\noindent
  \fcolorbox{violet}{bandviolet}{%
    \parbox{\dimexpr\linewidth-2\fboxsep-2\fboxrule}{%
      \vspace{2pt}%
      {\small\bfseries\color{violet}题干}\hspace{0.8em}#1%
      \vspace{2pt}}}%
  \par\medskip}

% v2: 保护题干与开头，长题可跨页，续页补上下文。
\newcommand{\qbreak}{\needspace{6\baselineskip}}
% 独立练习：ID、来源、题干、考点、答案标签（通常 ans:ID）。
\newcommand{\practice}[5]{%
  \qbreak\subhdr{练习 #1}\qstem{（#2）#3}%
  \textbf{考点}：#4\par
  详解见附录第~\pageref{#5}~页。\par\medskip}
% 附录开头：ID、来源、条件摘要；自动产生 ans:ID 标签。
\newcommand{\answerstart}[3]{%
  \qbreak\subhdr{#1\quad 详解}\label{ans:#1}\qstem{（#2）#3}}
\newcommand{\answercontinue}[2]{%
  \needspace{5\baselineskip}\subhdr{#1\quad 续解}%
  {\small\color{inkgray}本页接续：#2}\par\medskip}
% 延迟提供默认值，让旧讲义可在导言区保留自己的 \newcommand{\step}。
\AtBeginDocument{\providecommand{\step}[1]{\needspace{3\baselineskip}\par\medskip\textbf{#1}\par}}

% 前提旁注：只放在「需要理解的那一步」后面，给陈述 + 本题代入，不重证。
\newcommand{\aside}[1]{%
  \par\vspace{1pt}\noindent
  \hangindent=1.6em\hangafter=1
  {\small\color{inkgray}%
    \hspace*{0.6em}{\color{petal}$\triangleright$}\hspace{2pt}#1}%
  \par\vspace{2pt}}

%---------------------------------------------------------------- 表格
% 本机 MiKTeX 的 colortbl 与 array 版本不匹配（\rowcolor 会报错），
% 故表格统一用「上粗紫线 + 表头 + 中细线」的三线式，不依赖 colortbl。
\newcommand{\tabtop}{\par\vspace{3pt}{\color{violet}\hrule height 1.2pt}\vspace{1pt}}
\newcommand{\tabmid}{\par\vspace{1pt}{\color{petal}\hrule height 0.5pt}\vspace{1pt}}
\newcommand{\tabbot}{\par\vspace{3pt}{\color{rulepurple}\hrule height 0.5pt}\vspace{8pt}}
\newcommand{\thead}[1]{{\bfseries\color{violetdeep}#1}}

%---------------------------------------------------------------- 页眉页脚
\pagestyle{fancy}
\fancyhf{}
\renewcommand{\headrulewidth}{0.4pt}
\renewcommand{\headrule}{\hbox to\headwidth{\color{rulepurple}\leaders\hrule height \headrulewidth\hfill}}
\fancyhead[L]{\small\color{inkgray}\songti \doctag}
\fancyhead[R]{\small\color{inkgray}\songti \coursetag}
\fancyfoot[C]{\small\color{inkgray}\orn\enspace{\color{violet}\thepage}\enspace\orn}

% 每篇讲义在导言区设这两个值
\newcommand{\doctag}{}
\newcommand{\coursetag}{}
\newcommand{\settags}[2]{\renewcommand{\doctag}{#1}\renewcommand{\coursetag}{#2}}

%---------------------------------------------------------------- 数学记号
\newcommand{\vect}[1]{\boldsymbol{#1}}
\newcommand{\dd}{\,\mathrm{d}}
\newcommand{\ddx}{\frac{\mathrm{d}}{\mathrm{d}x}}
\addtolength{\jot}{2pt}



\usepackage{tikz}
\usetikzlibrary{arrows.meta,calc,positioning}
\usepackage{graphicx}
\setlength{\headheight}{22pt}
% theme 的 \tabmid 在 tabular 内会落进首格，改成 \noalign 形式
\renewcommand{\tabmid}{\noalign{\vskip 2pt{\color{petal}\hrule height 0.5pt}\vskip 3pt}}
\settags{Polarization Devices · L3 精讲}{光学（02）· 外教 Lecture 3}
% 节标题跟至少 4 行正文，避免标题孤零零落在页底
\let\subhdrtheme\subhdr
\renewcommand{\subhdr}[1]{\needspace{3\baselineskip}\subhdrtheme{#1}}
\let\hdrorig\hdr
\renewcommand{\hdr}[2]{\needspace{8\baselineskip}\hdrorig{#1}{#2}}
\setlength{\emergencystretch}{1.6em}
\tikzset{>={Stealth[length=2.0mm]},
  box/.style={draw=violet,rounded corners=2pt,fill=bandviolet,align=center,inner sep=3.5pt},
  boxc/.style={draw=petal,rounded corners=2pt,fill=bandcream,align=center,inner sep=3.5pt}}

\begin{document}

\doctitle{偏振器件与偏振光干涉（第三讲）}
{偏振棱镜 · 波片 · 补偿器 · 干涉颜色}
{依据：\texttt{Polarization Devices 1.pdf}（110 页，Lecture 3）\\
作业：课件 p60 的三题（题干见 §25，详解在附录 C）\\
光学（02）｜整理日期 2026-11-16}

%=====================================================================
\hdr{§0}{这一讲要解决什么问题}

上一讲把偏振算成了一个向量：两个正交分量加一个相位差。这一讲把那个向量做成器件，再用它做测量。
整堂课可以压成三句话：

\textbf{造一束纯的。}偏振片靠吸收，只能做到「大致偏振」；要靠双折射晶体把两束正交偏振光
\emph{在空间上分开}，才拿得到消光比 $10^{-5}$ 的纯光束。这是第一部分（棱镜）。

\textbf{改它的状态。}波片不吸收、不改变强度，只让两个分量错开一个相位；错开 $\pi/2$ 就把线偏振
变成圆偏振，错开 $\pi$ 就把振动面翻到轴对称的位置。补偿器是这个思想的连续可调版本。这是第二、三部分。

\textbf{看它的后果。}两个正交分量本来不干涉；一旦检偏器把它们投影到同一方向，干涉项就回来了，
于是延迟被翻译成颜色和条纹。这是第四部分，也是整堂课收口的地方——前面造的器件、做的波片，
最后都是为了让这件事发生。

{\small\color{inkgray}开篇地图：本讲把三个任务分别交给三件器件，最后在检偏器后面汇成干涉。}

\vspace{2pt}
\begin{center}
\begin{tikzpicture}[font=\footnotesize,x=1cm,y=1cm]
\node[box,text width=3.3cm] (q1) at (0,2.6) {\textbf{（1）要一束纯偏振光}\\ 偏振片吸收一半、玻璃堆不纯、单晶出两束};
\node[box,text width=3.3cm] (q2) at (0,0.2) {\textbf{（2）要改掉偏振态}\\ 线变圆、转动方位、换手性};
\node[box,text width=3.3cm] (q3) at (0,-2.2) {\textbf{（3）要量出未知延迟}\\ 应力、矿物、液晶盒};
\node[box,text width=3.7cm] (d1) at (5.6,2.6) {\textbf{偏振棱镜}（§2--§5）\\ 全反射扔掉 o 光，或让两束分开};
\node[box,text width=3.7cm] (d2) at (5.6,0.2) {\textbf{波片}（§8--§13）\\ 只改相位：$\Delta\varphi=2\pi\Delta n\,d/\lambda$};
\node[box,text width=3.7cm] (d3) at (5.6,-2.2) {\textbf{补偿器}（§14--§18）\\ 连续可变的已知延迟，与样品抵掉};
\node[box,text width=3.8cm] (r1) at (11.3,2.6) {纯度 $10^{-5}$ 的单束\\ 或两束正交偏振};
\node[box,text width=3.8cm] (r2) at (11.3,0.2) {$\lambda/4$ 造圆偏振；$\lambda/2$ 把方位翻到 $2\theta-\alpha$};
\node[box,text width=3.8cm] (r3) at (11.3,-2.2) {读出 $\Delta$（纳米）与 $\Delta n$，灵敏度 $\lambda/1000$};
\draw[->,violet] (q1)--(d1); \draw[->,violet] (q2)--(d2); \draw[->,violet] (q3)--(d3);
\draw[->,violet] (d1)--(r1); \draw[->,violet] (d2)--(r2); \draw[->,violet] (d3)--(r3);

\node[boxc,text width=9.0cm] (i1) at (5.6,-4.4) {\textbf{两条偏振光要干涉，必须先把它们投影到同一方向}：检偏器做的正是这件事（§19）};
\node[boxc,text width=9.0cm] (i2) at (5.6,-5.9) {$I=I_0\sin^2 2\varphi\;\sin^2\dfrac{\delta}{2}$：延迟被翻译成明暗与颜色（§20--§22）};
\node[boxc,text width=12.0cm] (a1) at (5.6,-7.3) {应用：DIC 显微镜 · 糖度计 · 椭偏仪 · Lyot 滤光片 · 光弹应力分析 · 电光调制（§7、§23）};
\draw[->,violet] (d2.south) -- ++(0,-0.55) -| (i1.north);
\draw[->,violet] (d3.south) -- ++(0,-0.35) -| (i1.north);
\draw[->,violet] (i1)--(i2); \draw[->,violet] (i2)--(a1);
\node[text width=15.4cm,align=left,font=\scriptsize,text=inkgray] at (5.6,-8.5)
{读法：上排三个任务各自对应一件器件（从左到右）；中下两框是把器件放进检偏器之后才出现的现象，它同时用到了「纯光束」和「已知延迟」。};
\end{tikzpicture}
\end{center}

%=====================================================================
\hdr{§0}{这一讲要解决什么问题}

上一讲把偏振算成了一个向量：两个正交分量加一个相位差。这一讲把那个向量做成器件，再用它做测量。
整堂课可以压成三句话：

\textbf{造一束纯的。}偏振片靠吸收，只能做到「大致偏振」；要靠双折射晶体把两束正交偏振光
\emph{在空间上分开}，才拿得到消光比 $10^{-5}$ 的纯光束。这是第一部分（棱镜）。

\textbf{改它的状态。}波片不吸收、不改变强度，只让两个分量错开一个相位；错开 $\pi/2$ 就把线偏振
变成圆偏振，错开 $\pi$ 就把振动面翻到轴对称的位置。补偿器是这个思想的连续可调版本。这是第二、三部分。

\textbf{看它的后果。}两个正交分量本来不干涉；一旦检偏器把它们投影到同一方向，干涉项就回来了，
于是延迟被翻译成颜色和条纹。这是第四部分，也是整堂课收口的地方——前面造的器件、做的波片，
最后都是为了让这件事发生。

\subhdr{0.1 页码路线：110 页课件对应到本讲义}

\begin{center}
\small
\tabtop
\begin{tabular}{@{}p{2.3cm}p{7.2cm}p{5.6cm}@{}}
\thead{课件页} & \thead{内容} & \thead{本讲义} \\
\tabmid
p2--p9 & 课程计划、偏振态回顾、Malus、Brewster、双折射、偏振史 & §1.2、§6 \\
\tabmid
p10--p20 & 偏振棱镜：共同原理、Nicol、Rochon、Wollaston、比较 & §2--§3、§5 \\
\tabmid
p21--p27 & 器件分类表：Glan 家族、薄膜分束立方 & §4 \\
\tabmid
p28--p48 & DIC、偏振测量、斯托克斯、零位法、糖度计、椭偏仪 & §6--§7 \\
\tabmid
p49--p52 & FAQ：损伤阈值、如何工作、如何选 & §5、§13 \\
\tabmid
p53--p67 & 波片：延迟、快慢轴、$\lambda/4$、$\lambda/2$、$\lambda$、色散、应用、FAQ & §8--§13 \\
\tabmid
p68--p83 & 补偿器：Babinet、Soleil--Babinet、Sénarmont、Berek、FAQ & §14--§18 \\
\tabmid
p84--p106 & 偏振光干涉：条件、公式、色偏振、干涉图、应用、FAQ & §19--§23 \\
\tabmid
p107--p109 & 三道例题（Nicol、波片、Babinet） & §24 \\
\tabmid
p110 & 课后作业三题 & §25、附录 B \\
\end{tabular}
\tabbot
\normalsize
\end{center}

\subhdr{0.2 记号（全篇统一）}

\begin{center}
\small
\tabtop
\begin{tabular}{@{}p{2.6cm}p{12.6cm}@{}}
\thead{记号} & \thead{含义} \\
\tabmid
$n_o$, $n_e$ & 寻常光、非常光的折射率；沿光轴传播时两者相等，都等于 $n_o$ \\
\tabmid
$\Delta n$ & $|n_e-n_o|$，双折射率的大小；方解石 $0.172$，石英 $0.009$ \\
\tabmid
$d$ & 晶片厚度（几何厚度） \\
\tabmid
$\Delta$, OPD & 光程差 $\Delta=\Delta n\,d$；相位差 $\delta=2\pi\Delta/\lambda$ \\
\tabmid
$\varphi$ & 波片快轴与起偏器透振方向的夹角（§20 的干涉公式用它） \\
\tabmid
$R$ & 样品或补偿器引入的延迟，单位常写成纳米 \\
\end{tabular}
\tabbot
\normalsize
\end{center}

%=====================================================================
\hdr{§1}{为什么要专门造一块棱镜}

\subhdr{1.1 三种「差不多」的办法都不够}

要一束纯的线偏振光，手边有三条路，但每条都差一点：

\textbf{偏振片（二向色性）。}它把不需要的那个分量\emph{吸收}掉，代价是白白损失一半能量，
而且吸收的能量变成热——强激光下会烧。它的消光比（想留的与漏过来的强度之比）大约 $10^{-3}$ 到 $10^{-4}$。

\textbf{玻璃堆。}以布儒斯特角叠一摞玻璃片，反射光越来越纯，但透射光永远只是「部分偏振」，
拿去当检偏器不够干净。

\textbf{单块双折射晶体。}光进去以后分成 o 光和 e 光，两束都在，出射是两种正交偏振的混合——
等于什么也没分出来。

\textbf{共同的缺口}：需要把两束正交偏振光\emph{在空间上分开}，并且只让其中一束出去。

\subhdr{1.2 双折射：两束光，两个折射率}

在方解石、石英这类各向异性晶体里，一束入射光会分成两束：o 光（ordinary）遵守普通的折射定律，
折射率恒为 $n_o$；e 光（extraordinary）一般\emph{不}遵守，它的折射率 $n_e(\theta)$ 随传播方向与光轴的夹角
$\theta$ 变化。只有沿着光轴传播时，两者才简并：

\[
n_e(0)=n_o .
\]

这就是罗雄棱镜第一块楔形「不分开」的原因（§3.1）：光沿着光轴走，晶体表现得像一个各向同性介质。

\begin{center}
\small
\tabtop
\begin{tabular}{@{}p{3.0cm}p{3.4cm}p{3.4cm}p{4.4cm}@{}}
\thead{晶体} & \thead{$n_o$} & \thead{$n_e$} & \thead{$\Delta n$ 与用途} \\
\tabmid
方解石（负晶体） & 1.658 & 1.486 & $0.172$，天然双折射最强 $\to$ 做\emph{分束}棱镜 \\
\tabmid
石英（正晶体） & 1.544 & 1.553 & $0.009$，小且均匀 $\to$ 做\emph{波片}与补偿器 \\
\end{tabular}
\tabbot
\normalsize
\end{center}

\textbf{一个数量级就能说明为什么分工不同。}取 $1\ \mathrm{cm}$ 厚的方解石，光程差
$\Delta=\Delta n\,d=0.172\times10^{-2}\ \mathrm{m}=1.72\ \mathrm{mm}$，在 $589\ \mathrm{nm}$ 处是
$\dfrac{1.72\times10^{-3}}{589\times10^{-9}}\approx 2920$ 个波长。要让这样一块晶体当四分之一波片，
你得把厚度控制到 $2920$ 分之一以内——做不到。反过来，石英 $\Delta n$ 小，做波片时厚度对误差不敏感，
但拿它做分束棱镜，偏折角小得没法用。

\subhdr{1.3 两种分离机制}

课件把偏振棱镜分成两族，分法是「两束光最后去了哪里」：

\textbf{全反射型（Nicol、Glan 家族）。}在晶体中间夹一层\emph{折射率介于 $n_e$ 与 $n_o$ 之间}的介质
（经典做法是加拿大树胶，$n=1.55$）。对 o 光，界面是「光密到光疏」，超过临界角就全反射，被侧壁吸收；
对 e 光，$n_e<n_{\text{胶}}$，界面反而是「光疏到光密」，全反射在物理上不可能发生，于是直穿过去。
一块器件只出一个光束，纯度极高。

\textbf{偏折型（Rochon、Wollaston）。}不做全反射，而是让两束光在界面处\emph{偏折角不同}，
出去时分开一个小角度。两束都能用，或者一束直、一束偏。

%=====================================================================
\hdr{§2}{尼科尔棱镜：用全反射把 o 光扔掉}

\subhdr{2.1 构造}

两块方解石菱面体磨成端面与光轴约成 $68^\circ$ 的楔形，用加拿大树胶（$n_b=1.55$）粘起来，
侧面涂黑。光从端面射入，o 光在胶层上全反射，打到涂黑的侧壁被吸收；e 光几乎不偏折地穿出。
出射光在\emph{主截面}内振动，位置横向平移 $0.05$--$0.2\ \mathrm{mm}$（$25\ \mathrm{mm}$ 的棱镜约 $0.1\ \mathrm{mm}$），
方向与入射光平行。

\subhdr{2.2 临界角：$68^\circ$ 这个数字怎么来的}

全反射的门槛是临界角。对 o 光，从方解石（$n_o$）射向树胶（$n_b$）：

\[
\sin\theta_c=\frac{n_b}{n_o}=\frac{1.55}{1.658}=0.9349
\quad\Longrightarrow\quad
\theta_c=69.2^\circ .
\]

也就是说，o 光只要以大于 $69.2^\circ$ 的角度打到胶层上，就必定全反射。棱镜的切割让 o 光在胶层上的入射角
稳定地大于这个值，所以它「总是」被扔掉。课件写 $68^\circ$，指的是端面与光轴的夹角（磨制时的几何量），
与胶层上的入射角不是同一个角，两者差着一个由棱镜几何决定的量；判断准则始终是上面那个临界角。

\textbf{e 光为什么不会跟着反射？}因为 $n_e=1.486<n_b=1.55$，e 光是从\emph{光疏}射向\emph{光密}，
无论入射角多大都不可能全反射（临界角只在「密到疏」时存在）。这件事与角度无关，是材料决定的——
所以尼科尔棱镜对 e 光的透过是「内生」的，不靠微调几何。

\subhdr{2.3 它为什么被淘汰，又留下了什么}

尼科尔棱镜的视场只有大约 $\pm14^\circ$，口径做不大，还要消耗大量优质方解石（天然大晶体稀缺）。
今天实验室里见到的是它的直系后代：把树胶换成空气隙的 Glan--Foucault（耐强光），
或把几何改好的 Glan--Taylor（视场大、阈值高）。原理一字未变：\emph{让 o 光在折射率夹层上全反射}。

\begin{center}
\begin{tikzpicture}[font=\footnotesize,x=1cm,y=1cm]
% 棱镜外形（第一块：斜线以上；第二块：斜线以下的三角）
\fill[bandviolet] (0,0) -- (0.8,0) -- (4.8,1.45) -- (4.8,2.0) -- (0,2.0) -- cycle;
\fill[bandcream] (0.8,0) -- (4.8,0) -- (4.8,1.45) -- cycle;
\draw[thick] (0,0) rectangle (4.8,2.0);
\draw[violet,thick] (0.8,0) -- (4.8,1.45);
\draw[line width=2.6pt] (4.0,0) -- (4.8,0);
\node[font=\scriptsize,text=violet,rotate=21] at (1.30,0.32) {树胶 $n_b=1.55$};
% 入射
\draw[->,thick,orange!80!black] (-1.2,1.0) -- (3.5,1.0);
\node[font=\scriptsize,color=orange!80!black] at (-0.55,1.22) {自然光};
% o 光反射到底面
\draw[->,thick,densely dashed,blue!70!black] (3.5,1.0) -- (4.68,0.06);
\node[font=\scriptsize,color=blue!70!black] at (3.62,0.50) {o 光};
% e 光穿出
\draw[->,thick,red!70!black] (3.5,1.0) -- (6.1,0.94);
\node[font=\scriptsize,color=red!70!black] at (5.45,1.18) {e 光};
\end{tikzpicture}

{\small\color{inkgray}图 1\quad 尼科尔棱镜的示意图：斜线是树胶层（$n_b=1.55$，介于 $n_o=1.658$ 与
$n_e=1.486$ 之间）。o 光在胶层上的入射角约 $70^\circ$，超过临界角 $69.2^\circ$，被全反射到底面
（涂黑）吸收；e 光从光疏射向光密，不可能全反射，几乎不偏折地穿出。}
\end{center}

%=====================================================================
\hdr{§3}{罗雄与沃拉斯顿：让两束都出去}

\subhdr{3.1 罗雄棱镜：第一块为什么不分开}

罗雄棱镜由两块楔形方解石胶合。\textbf{前一块的光轴平行于传播方向}，后一块的光轴与它垂直。
在第一块里，光沿光轴传播，两个分量看到的折射率都是 $n_o$，晶体表现得像各向同性介质，
\emph{根本不分裂}；到了第二块，光轴垂直于传播方向，两个分量才分别看到 $n_o$ 与 $n_e$，
在界面上按不同折射率偏折。

结果是一束直、一束偏：e 光（或按切割方式而定的一束）沿原方向出射，另一束偏出 $1^\circ$--$10^\circ$。
由于没有用折射率匹配的胶，工作带宽可达 $350\ \mathrm{nm}$ 到 $2.5\ \mu\mathrm{m}$，也耐强光；
缺点是偏向的那束色散明显（偏向角跟着 $\Delta n(\lambda)$ 变）。

\subhdr{3.2 沃拉斯顿棱镜：两束对称分开}

沃拉斯顿也是两块直角楔形方解石胶合，但\textbf{两块的光轴互相垂直，且都躺在端面内}。
光在第一块里分成 o、e 两个分量；到了第二块，由于光轴转了 $90^\circ$，\emph{原来的o光变成了e光，
原来的e光变成了o光}——角色互换。于是两个分量在界面上朝着相反方向偏折：

\[
\varepsilon=2\alpha\,\Delta n ,\qquad
\text{两束夹角}=2\varepsilon=4\alpha\,\Delta n\quad(\text{小楔角近似}).
\]

课件给的数量级与此一致：$\Delta n=0.172$、$\alpha=5^\circ$ 时，罗雄偏 $\approx1.7^\circ$，
沃拉斯顿两束分开 $\approx3.4^\circ$。（这一步用到的近似是「楔角小、偏向角与厚度无关」，
精确的偏向角还要乘上表面折射的放大因子，量级不变。）

\subhdr{3.3 两者的工程差别}

\begin{center}
\small
\tabtop
\begin{tabular}{@{}p{3.6cm}p{5.6cm}p{5.6cm}@{}}
\thead{} & \thead{罗雄 Rochon} & \thead{沃拉斯顿 Wollaston} \\
\tabmid
出射 & 一束直、一束偏 & 两束都偏，对称各 $\pm\varepsilon$ \\
\tabmid
输出数 & 通常只用直的那束 & 两束都能用，强度可相等 \\
\tabmid
色散 & 偏向束色散大 & 分束角随 $\lambda$ 变 \\
\tabmid
典型用途 & 激光系统、需要准直的场合 & 剪切干涉、DIC 显微镜、偏振复用 \\
\end{tabular}
\tabbot
\normalsize
\end{center}

\textbf{为什么 DIC 显微镜偏爱沃拉斯顿？}因为 DIC 需要两束\emph{等强度、对称剪切}的光穿过样品，
再由第二块沃拉斯顿精确合回；对称几何保证两束振幅相等、公共相位抵消，才能把样品的相位梯度变成
「一边亮一边暗」的浮雕感。罗雄一束强一束弱，做不出这种对称的对比。

\begin{center}
\begin{tikzpicture}[font=\footnotesize,x=1cm,y=1cm]
% (a) 罗雄
\fill[bandviolet] (0,0) -- (1.6,0) -- (4.4,1.8) -- (0,1.8) -- cycle;
\fill[bandcream] (1.6,0) -- (4.4,0) -- (4.4,1.8) -- cycle;
\draw[thick] (0,0) rectangle (4.4,1.8);
\draw[violet,thick] (1.6,0) -- (4.4,1.8);
\draw[<->,violet,thick] (0.45,0.95) -- (0.45,1.50);
\node[font=\scriptsize,text=violet] at (0.45,1.78) {轴 $\parallel$ 光};
\draw[<->,violet,thick] (3.85,0.60) -- (3.85,1.10);
\node[font=\scriptsize,text=violet] at (3.85,0.35) {轴 $\perp$ 光};
\draw[->,thick] (-1.0,0.9) -- (3.1,0.9);
\draw[->,thick,red!70!black] (3.1,0.9) -- (6.2,0.9);
\draw[->,thick,blue!70!black] (3.1,0.9) -- (6.2,0.60);
\node[font=\scriptsize,color=red!70!black] at (5.55,1.06) {直：e 光};
\node[font=\scriptsize,color=blue!70!black] at (5.55,0.40) {偏：o 光};
\node[font=\scriptsize] at (2.2,-0.38) {(a) 罗雄：一束直、一束偏};

% (b) 沃拉斯顿
\begin{scope}[yshift=-3.4cm]
\fill[bandviolet] (0,0) -- (1.6,1.8) -- (0,1.8) -- cycle;
\fill[bandcream] (0,0) -- (4.4,0) -- (4.4,1.8) -- (1.6,1.8) -- cycle;
\draw[thick] (0,0) rectangle (4.4,1.8);
\draw[violet,thick] (1.6,1.8) -- (4.4,0);
\draw[<->,violet,thick] (0.55,0.70) -- (0.55,1.25);
\node[font=\scriptsize,text=violet] at (0.55,1.52) {轴 1};
\draw[<->,violet,thick] (3.15,0.95) -- (3.85,0.95);
\node[font=\scriptsize,text=violet] at (3.50,1.18) {轴 2};
\draw[->,thick] (-1.0,0.9) -- (2.95,0.9);
\draw[->,thick,red!70!black] (2.95,0.9) -- (6.2,1.42);
\draw[->,thick,blue!70!black] (2.95,0.9) -- (6.2,0.38);
\node[font=\scriptsize,color=red!70!black] at (5.55,1.60) {$+\varepsilon$};
\node[font=\scriptsize,color=blue!70!black] at (5.55,0.20) {$-\varepsilon$};
\node[font=\scriptsize] at (2.2,-0.38) {(b) 沃拉斯顿：两束对称分开};
\end{scope}
\end{tikzpicture}

{\small\color{inkgray}图 2\quad 罗雄与沃拉斯顿的几何对比（示意）。罗雄的第一块里光沿轴（轴 1）走、
不分裂，第二块（轴 2）才偏折，所以一束直、一束偏；沃拉斯顿的两块轴互相垂直（轴 1、轴 2
都在端面内），界面处 o/e 角色互换，两束朝相反方向偏。}
\end{center}

%=====================================================================
\hdr{§4}{格兰家族与薄膜分束立方}

\subhdr{4.1 同一原理的三种改良}

\textbf{Glan--Thompson}：两块方解石用加拿大树胶粘合，光轴垂直于入射面。视场大、透过率高，
但树胶是有机物，功率一高就烧——只能在低功率下用。

\textbf{Glan--Foucault}：把树胶换成\emph{空气隙}。没有有机物，耐功率高、能进紫外，代价是接受角变小。

\textbf{Glan--Taylor}：空气隙加改良几何，兼得大视场与高阈值，是现代激光实验室的标准偏振器。

三者与尼科尔的关系是「同一个全反射思路 + 更好的胶层」。判断选型时记住一句话：
\textbf{胶层是薄弱环节}，功率密度一上去就换成空气隙。

\subhdr{4.2 薄膜分束立方：不是双折射}

还有一类器件长得像偏振棱镜，原理却完全不同：两块玻璃直角棱镜之间镀一层介质膜，
膜让 p 分量透射、s 分量反射。它靠\emph{镀膜}而不是晶体的双折射，所以：
优点是没有晶体尺寸限制、可做大口径、带宽可以设计；缺点是消光比通常不如晶体棱镜，
而且 s/p 分离靠入射角，斜入射时性能下降。

\textbf{一句话区分}：晶体棱镜利用「两个分量在晶体里折射率不同」；薄膜立方利用「两个分量在膜堆里
反射率不同」。

\subhdr{4.3 常见器件一览}

\begin{center}
\small
\tabtop
\begin{tabular}{@{}p{3.5cm}p{4.6cm}p{2.0cm}p{4.6cm}@{}}
\thead{器件} & \thead{机制} & \thead{输出} & \thead{特点} \\
\tabmid
Nicol & 树胶层全反射 & 单束 & 历史器件，视场小，已淘汰 \\
\tabmid
Glan--Thompson & 树胶层全反射 & 单束 & 视场大，低功率 \\
\tabmid
Glan--Foucault & 空气隙全反射 & 单束 & 高阈值、紫外，接受角小 \\
\tabmid
Rochon & 偏折（第一块沿轴） & 双束，一束直 & 宽带、耐强光 \\
\tabmid
Wollaston & 偏折（轴互换） & 双束对称 & 剪切、DIC、复用 \\
\tabmid
薄膜立方 & 介质膜反射/透射 & 双束 & 大口径，消光比一般 \\
\end{tabular}
\tabbot
\normalsize
\end{center}

%=====================================================================
\hdr{§5}{怎么挑一块棱镜：读规格表的顺序}

课件给了一份很实用的检查顺序，按重要性排：

\textbf{（1）波长范围}决定材料。方解石到 $2.3\ \mu\mathrm{m}$；近红外用 YVO$_4$；紫外用 MgF$_2$ 或 $\alpha$-BBO。
方解石在 $230\ \mathrm{nm}$ 以下强烈吸收，这是硬限制，不是镀膜能补救的。

\textbf{（2）消光比}。常规 $10^{-5}$；精密偏振测量要 $10^{-6}$ 或更好。注意：增透膜能提高\emph{透过率}，
但提高不了消光比——消光比由晶体质量、胶层和杂散光决定。

\textbf{（3）接受角}。必须覆盖光束的全部锥角，包括显微物镜的会聚锥。这一条最容易被忽略。

\textbf{（4）损伤阈值}。脉冲激光用 $\mathrm{J/cm^2}$，连续光用 $\mathrm{W/cm^2}$。胶合件是最弱的一环。

\textbf{（5）透过率与面型}。

\warnbox{\textbf{最常见的选型错误}：买了一块消光比极好、但接受角很窄的棱镜，然后把它放进会聚光束里。
光束一倾斜，o 光在胶层上的入射角降回临界角以下，本该被反射的光漏了过去，纯度立刻崩塌。
\textbf{消光比是在准直、正入射条件下标的}，用之前先确认自己的光是不是准直的。}

%=====================================================================
\hdr{§6}{偏振怎么测：三个层次}

\subhdr{6.1 转动检偏器：能测三个斯托克斯参数}

把待测光过一个可旋转的线偏振片，后面接探测器。设检偏器转到角度 $\theta$，透过的强度由 Malus 定律给出：

\[
I(\theta)=\tfrac12\left(S_0+S_1\cos2\theta+S_2\sin2\theta\right).
\]

这里的 $S_0$ 是总强度，$S_1=I(0^\circ)-I(90^\circ)$ 是水平与垂直线偏振的差，$S_2=I(45^\circ)-I(135^\circ)$
是 $\pm45^\circ$ 线偏振的差。\emph{为什么是这个形式}：$\cos2\theta$、$\sin2\theta$ 恰好是「转过 $\theta$ 的检偏器
看到水平分量与 $45^\circ$ 分量的权重」，把 $I(\theta)$ 做傅里叶分解，前三个系数就是 $S_0,S_1,S_2$。
检验：代入 $\theta=0$，右端 $\tfrac12(S_0+S_1)=I(0^\circ)$，核对成立；代入 $45^\circ$，$\tfrac12(S_0+S_2)=I(45^\circ)$，核对成立。

\textbf{它测不到 $S_3$。}因为线偏振片对左旋、右旋圆偏振光的响应完全一样（都透过一半），
手性信息在投影时丢掉了。补救办法：在检偏器前加一片四分之一波片，把圆偏振先变成线偏振，
转动检偏器就能读出 $S_3$——这就是完整的斯托克斯偏振测量。

\subhdr{6.2 零位法：为什么「找最暗」比「读最亮」准}

直接比较两个亮场的强度，人眼和探测器的线性度都有限；但在接近全暗的地方，强度的相对变化被放大，
判断「两边一样暗」比判断「两边一样亮」精确得多。这就是零位法（null method）的核心。

\textbf{半荫法}是它的工程版本：视场被分成两半，检偏器转到某个位置时两半\emph{一样暗}。
这个位置对光源强度起伏不敏感，也不需要读绝对强度，所以精度远高于直接测极小值。
糖度计（§7.2）用的就是这一套。

\subhdr{6.3 不确定度从哪里来}

三个主要来源：偏振器有限的消光比（暗场不可能绝对黑，给最小值设了地板）、机械转角读数、
以及温度与装夹应力引起的寄生双折射。最后一项常被忽略：棱镜夹得太紧，玻璃里的应力本身就会
产生几个纳米的延迟。

%=====================================================================
\hdr{§7}{三个工程案例：器件怎么用在系统里}

\subhdr{7.1 微分干涉相衬（DIC）显微镜}

活细胞几乎透明，明场下看不见。DIC 的做法是：聚光镜里的沃拉斯顿棱镜把照明光分成两束正交偏振、
横向错开不到 $1\ \mu\mathrm{m}$ 的光；两束光穿过样品时走过略微不同的路径，积累一个与
\emph{折射率梯度}成正比的微小光程差；物镜后的第二块沃拉斯顿把它们重新合到一起；
最后检偏器把两个正交分量投影到同一方向，让它们干涉。于是「相位梯度」变成了「明暗对比」，
细胞边缘出现典型的浮雕感。

\textbf{与本讲的联系}：没有棱镜就没有两束正交光，没有检偏器就没有干涉——这正是 §19 的
Fresnel--Arago 结论在显微镜里的实现。

\subhdr{7.2 糖度计：把旋光角读成浓度}

蔗糖等手性分子会让线偏振光的振动面旋转一个角度 $\alpha$，且

\[
\alpha=[\alpha]_D\,c\,L ,
\]

其中 $[\alpha]_D$ 是比旋光（钠 D 线，$589\ \mathrm{nm}$），$c$ 是浓度，$L$ 是光程。
蔗糖的比旋光约为 $+66.5^\circ$（每 $\mathrm{dm}$、每 $\mathrm{g/mL}$ 的量级），
葡萄糖 $+52.7^\circ$，果糖 $-92.4^\circ$——正负号就是「右旋/左旋」。

测量步骤就是零位法：先放纯水，把检偏器调到两半等暗（零点）；换样品，振动面被转过 $\alpha$，
再把检偏器转回来，转过的角度就是 $\alpha$。仪器刻度直接按 $^\circ\mathrm{Z}$（糖度）标定，
所以读数是浓度而不是角度。

\textbf{为什么必须固定波长}：旋光角随波长变（旋光色散），换波长刻度就全错。

\subhdr{7.3 椭偏仪：把偏振态的变化变成膜厚}

线偏振光斜入射到表面，p 分量与 s 分量的反射系数不同（振幅比与相位都不同），
反射光一般变成椭圆偏振。椭偏仪测的正是这两个量：

\[
\tan\Psi=\frac{|r_p|}{|r_s|},\qquad \Delta=\delta_p-\delta_s,\qquad
\rho=\frac{r_p}{r_s}=\tan\Psi\,e^{i\Delta}.
\]

$\Psi$ 与 $\Delta$ 由转动检偏器（或转动补偿器）测出，再反演膜厚与折射率。
精度能到 $0.01\ \mathrm{nm}$，是非破坏测量薄膜的标准手段。

%=====================================================================
\hdr{§8}{波片：只改相位，不改强度}

\subhdr{8.1 一块晶片就是两台速度不同的电梯}

把双折射晶体切成平行薄片，\emph{让光轴躺在片子里}，光正入射、沿垂直光轴的方向走。这时两个正交分量
分别沿「快轴」和「慢轴」振动，各走各的速度，走完厚度 $d$ 之后：

\[
\Delta=\Delta n\,d,\qquad \delta=\frac{2\pi}{\lambda}\Delta n\,d .
\]

$\Delta$ 是光程差（单位纳米），$\delta$ 是相位差（单位弧度）。这是全篇最重要的一个式子：\emph{厚度把相位差
变成可设计的量}。

\textbf{小例。}$1\ \mathrm{mm}$ 厚的石英（$\Delta n=0.009$）：$\Delta=9\ \mu\mathrm{m}$，在 $589\ \mathrm{nm}$
处相当于 $15.3$ 个波长的延迟。注意这个「$15.3$」——整数部分没有物理后果（相当于 15 个整波，
两个分量回到原来的相对相位），只有小数部分 $0.3$ 起作用。这就是为什么实用的波片要「零级」：
把厚度做到最小，让小数部分干净可控。

\subhdr{8.2 快轴、慢轴与正负晶体}

\textbf{快轴}是折射率小的方向，光在这里跑得快、相位超前；\textbf{慢轴}正交，折射率大、相位落后。

\begin{center}
\small
\tabtop
\begin{tabular}{@{}p{3.2cm}p{2.4cm}p{2.4cm}p{6.6cm}@{}}
\thead{晶体} & \thead{$n_e$ 与 $n_o$} & \thead{类型} & \thead{快轴是哪条} \\
\tabmid
方解石 & $n_e<n_o$ & 负晶体 & e 轴（光轴方向）\\
\tabmid
石英、MgF$_2$ & $n_e>n_o$ & 正晶体 & o 轴（垂直光轴方向）\\
\end{tabular}
\tabbot
\normalsize
\end{center}

\textbf{为什么这条要单独记}：写琼斯矢量、判断旋向、判断补偿器读数的正负，全都要知道「哪条是快轴」。
厂商会在镜框上刻一条线或打一个点标快轴；没有标记时，把波片放在正交偏振器之间转 $45^\circ$ 看颜色，
与已知波片比对即可定出。

\subhdr{8.3 它与偏振片的根本区别}

偏振片\emph{吸收}一个分量，强度必然下降；波片\emph{两个分量都保留}，只改它们的相对相位，
理想情况下不吸收任何能量。这条差别决定了用途：波片可以放在激光谐振腔里、放在高功率光路里，
偏振片不行。

\textbf{它也不能造出偏振光。}自然光的两个分量互不相干，波片把它们错开相位之后仍然互不相干，
出来还是自然光。要造偏振光，必须先有偏振片（或棱镜、或布儒斯特面）。

\subhdr{8.4 零级与多级}

\begin{center}
\small
\tabtop
\begin{tabular}{@{}p{3.0cm}p{4.6cm}p{7.0cm}@{}}
\thead{} & \thead{零级（$m=0$）} & \thead{多级} \\
\tabmid
厚度 & $d=\lambda/(4\Delta n)$，最薄 & $d=(m+\tfrac14)\lambda/\Delta n$ \\
\tabmid
优点 & 对波长、温度不敏感 & 好磨、成本低、结实 \\
\tabmid
缺点 & 太薄，难加工 & 延迟随 $\lambda$ 与温度漂移大 \\
\tabmid
数值（石英 $589\ \mathrm{nm}$） & $16.2\ \mu\mathrm{m}$ & 20 级约 $1.31\ \mathrm{mm}$ \\
\end{tabular}
\tabbot
\normalsize
\end{center}

%=====================================================================
\hdr{§9}{四分之一波片：把线偏振拧成圆}

\subhdr{9.1 厚度条件}

要求 $\delta=\pi/2$，即光程差为四分之一个波长：

\[
\Delta n\,d=\frac{\lambda}{4}\quad\Longrightarrow\quad d_{\lambda/4}=\frac{\lambda}{4\Delta n}.
\]

\subhdr{9.2 为什么必须成 $45^\circ$：一个完整的推导}

取快轴为 $x$、慢轴为 $y$。入射是线偏振，振幅 $E_0$，与快轴成 $\varphi$ 角，写成两个分量：

\[
\vect E_{\text{in}}=E_0\begin{pmatrix}\cos\varphi\\ \sin\varphi\end{pmatrix}.
\]

过片之后慢轴分量多落后 $\delta$（写成 $e^{-i\delta}$ 的因子）：

\[
\vect E_{\text{out}}=E_0\begin{pmatrix}\cos\varphi\\ \sin\varphi\,e^{-i\delta}\end{pmatrix}
=E_0\begin{pmatrix}\cos\varphi\\ -i\sin\varphi\end{pmatrix}\quad(\delta=\pi/2).
\]

要让端点走一个\emph{圆}，需要两个条件同时成立：两个分量的振幅相等（$|\cos\varphi|=|\sin\varphi|$，
即 $\varphi=45^\circ$），并且相位差正好 $90^\circ$。取 $\varphi=45^\circ$：

\[
\vect E_{\text{out}}=\frac{E_0}{\sqrt2}\begin{pmatrix}1\\ -i\end{pmatrix},
\]

两个分量等幅、相差 $90^\circ$，端点做圆周运动——圆偏振。

\textbf{三个反面情形同样重要}：$\varphi=0^\circ$ 或 $90^\circ$ 时只有一个分量，波片无事可做，出来还是线偏振；
其他角度两个分量不等幅，出来是椭圆偏振。

\begin{center}
\small
\tabtop
\begin{tabular}{@{}p{5.6cm}p{3.4cm}p{5.0cm}@{}}
\thead{入射态（对快轴）} & \thead{出射态} & \thead{备注} \\
\tabmid
线偏振，平行于某一轴 & 线偏振，不变 & 只有一个分量 \\
\tabmid
线偏振，$45^\circ$ & 圆偏振 & 振幅相等 + 相差 $90^\circ$ \\
\tabmid
线偏振，其他角度 & 椭圆偏振 & 一般情形 \\
\tabmid
圆偏振 & 线偏振 & 反向操作，方位由片子取向定 \\
\tabmid
椭圆偏振，轴对齐 & 椭圆偏振，转过 $90^\circ$ & 长短轴互换 \\
\tabmid
自然光 & 自然光 & 两个分量不相干，相位差无意义 \\
\end{tabular}
\tabbot
\normalsize
\end{center}

\subhdr{9.3 一个实用的判别法}

圆偏振光过检偏器时，\emph{无论检偏器转到哪个角度，透过的强度都是 $I_0/2$}——因为任何方向上的投影
平均下来都是一半。这与自然光的表现相同，单靠检偏器分不开两者。加一片四分之一波片就能分开：
圆偏振会被它变回线偏振，转动检偏器就出现消光；自然光仍然是自然光，怎么转都是 $I_0/2$。

%=====================================================================
\hdr{§10}{二分之一波片：把振动面翻到轴对称的位置}

\subhdr{10.1 $\delta=\pi$ 意味着「一个分量反号」}

半个波长的延迟等价于把慢轴分量乘以 $-1$：$e^{-i\pi}=-1$。两个分量一正一负，
合成矢量的方向就是原来振动方向关于快轴的\emph{镜像}。

\subhdr{10.2 公式 $\alpha_{\text{out}}=2\theta-\alpha_{\text{in}}$ 的来历}

设快轴在实验室坐标里成 $\theta$ 角，入射振动方向成 $\alpha_{\text{in}}$ 角。
换到「以快轴为 $0^\circ$」的坐标系里，入射角是 $\alpha_{\text{in}}-\theta$；反号之后出射角是
$-(\alpha_{\text{in}}-\theta)$；再换回实验室坐标：

\[
\alpha_{\text{out}}=\theta-\left(\alpha_{\text{in}}-\theta\right)=2\theta-\alpha_{\text{in}} .
\]

\textbf{读法}：出射方向是入射方向关于快轴的镜像；把波片转 $\theta$，出射方向转 $2\theta$——
所以半波片被叫作「杠杆比为 2 的偏振旋转器」。

\textbf{两个特例}值得记住：$\theta=0$（快轴与入射方向重合）时 $\alpha_{\text{out}}=-\alpha_{\text{in}}$，
即关于快轴镜像；$\theta=45^\circ$ 时 $\alpha_{\text{out}}=90^\circ-\alpha_{\text{in}}$，
若入射也是 $45^\circ$，出射就是 $90^\circ$（正交）。

\subhdr{10.3 圆偏振进来会换手}

左旋进、右旋出。原因很直接：半波片把慢轴分量反号，等于把电场旋转的方向反过来。
这一条在 3D 眼镜、光隔离器里都要用到。

\subhdr{10.4 为什么它比「两片偏振片」好用}

两片偏振片也能「转动」偏振方向：转第一片即可。但那要吸收掉一半以上的光。半波片只改相位、
几乎不吸收，还能保持光束方向不变。所以在激光腔内、在加工头的偏振切换、在光纤偏振控制器里，
用半波片而不是偏振片。

%=====================================================================
\hdr{§11}{全波片：在单色光里「看起来不存在」}

$\delta=2\pi$ 时两个分量回到同相，偏振态完全复原，所以全波片在单色光下\emph{看不出任何效果}。
它的用处出现在白光里：设计波长满足 $\Delta=\lambda$，其他波长不满足，各自带着不同的残余延迟出来，
于是视场出现颜色（§21 的色偏振）。

三种波片的厚度比恒为

\[
d_{\lambda/4}:d_{\lambda/2}:d_{\lambda}=1:2:4 .
\]

%=====================================================================
\hdr{§12}{五个必须会预测的操作和三条规则}

\begin{center}
\small
\tabtop
\begin{tabular}{@{}p{4.4cm}p{5.8cm}p{5.2cm}@{}}
\thead{想要} & \thead{用什么} & \thead{关键条件} \\
\tabmid
自然光 $\to$ 线偏振 & 偏振片（棱镜或二向色片） & 波片做不到 \\
\tabmid
线偏振 $\to$ 圆偏振 & $\lambda/4$ 片 & 快轴与入射成 $45^\circ$ \\
\tabmid
圆偏振 $\to$ 线偏振 & $\lambda/4$ 片 & 出射方位随片子取向转 \\
\tabmid
线偏振 $\to$ 另一方位 & $\lambda/2$ 片 & 转到 $\theta=(\alpha_{\text{in}}+\alpha_{\text{out}})/2$ \\
\tabmid
任意态 $\to$ 任意态 & 可转 $\lambda/2$ + 可转 $\lambda/4$ & 光纤偏振控制器的标准配置 \\
\end{tabular}
\tabbot
\normalsize
\end{center}

\textbf{三条规则}，做题时能省一半时间：

\begin{enumerate}[label=（\arabic*）,leftmargin=2.4em]
\item 波片的轴与入射振动方向平行或垂直时，什么都不会发生（只有一个分量）。
\item 延迟可以相加：两片 $\lambda/4$ 轴交叉叠起来等于 $\lambda/2$；两片 $\lambda/2$ 同向叠加等于全波。
\item 双程加倍：$\lambda/4$ 片后面放一面镜子，反射光再穿一次，等于 $\lambda/2$。
这就是反射式液晶屏和双程隔离器的原理。
\end{enumerate}

%=====================================================================
\hdr{§13}{色散、消色差与几个常见坑}

\subhdr{13.1 一片波片只在设计波长上是它自己}

延迟 $\delta=2\pi\Delta n\,d/\lambda$ 里，$\Delta n$ 与 $\lambda$ 都随波长变，所以「$633\ \mathrm{nm}$ 的
半波片」拿到 $532\ \mathrm{nm}$ 用就不再是半波片。算一下：石英半波片
$d=633/(2\times0.0091)=34.8\ \mu\mathrm{m}$，换到 $532\ \mathrm{nm}$，

\[
\delta=2\pi\frac{0.0091\times34.8\ \mu\mathrm{m}}{532\ \mathrm{nm}}
=2\pi\times\frac{316.5}{532}=2\pi\times0.595 ,
\]

即 $214^\circ$，不是 $180^\circ$——差出去 $34^\circ$，圆偏振会变成明显的椭圆。

\subhdr{13.2 三条消色差路线}

\textbf{两片不同材料、轴交叉}：石英与 MgF$_2$ 的色散一正一负，互相抵消，在约一个倍频程内
延迟近乎恒定。

\textbf{Pancharatnam（超消色差）}：三片以上、轴依次取 $0^\circ,60^\circ,0^\circ$，在中心波长
$\pm30\%$ 内平坦。天文偏振测量常用。

\textbf{Fresnel 菱形}：不用双折射，用玻璃内部两次全反射的相位差，宽带且便宜，代价是光束横向平移。

\subhdr{13.3 四个常被问到的点}

\textbf{「波片会减慢光速吗？」}会——慢轴分量确实走得慢，但两个分量走的是同样的几何厚度，
差的只是相位，总强度不变。

\textbf{「零级为什么比多级稳？」}因为延迟正比于厚度；厚度大 $20$ 倍，同样的温度或波长漂移
就被放大 $20$ 倍。

\textbf{「能拿波片当偏振片吗？」}不能。它没有吸收轴，只改相位；造偏振光必须另有偏振器。

\textbf{「斜着用行不行？」}斜入射会改变有效延迟与光程，零级片最宽容，会聚光束里最好别用。

%=====================================================================
\hdr{§14}{补偿器：把延迟做成可调的旋钮}

\subhdr{14.1 固定波片解决不了的问题}

一片 $\lambda/4$ 片只会给你 $\pi/2$，一个数。可真实样品——受力的塑料件、矿物薄片、液晶盒、生物纤维——
的双折射是\emph{未知}的，延迟可能是 $37\ \mathrm{nm}$，也可能是 $1.2\ \mu\mathrm{m}$。
你需要的不是一片固定波片，而是一个\emph{已知、连续可变}的延迟，用它去把样品的延迟抵消掉。

\subhdr{14.2 差动楔形：用厚度差造延迟}

两块光轴互相垂直的晶体楔，厚度分别为 $t_1$、$t_2$。由于轴垂直，一块的快轴正是另一块的慢轴，
两者的延迟\emph{相减}：

\[
\Delta=\Delta n\,(t_1-t_2),\qquad
\delta=\frac{2\pi}{\lambda}\Delta n\,(t_1-t_2).
\]

只要 $t_1=t_2$，净延迟就是零；改变相对厚度，延迟就从零连续变到几个波长。
这就是所有补偿器的共同骨架，区别只在「怎么改变 $t_1-t_2$」。

%=====================================================================
\hdr{§15}{巴比涅补偿器：平移楔形，看黑带移动}

\subhdr{15.1 构造与读数}

两片石英楔，楔角 $2^\circ$--$5^\circ$，光轴互相垂直，斜面相对放置；一片固定，另一片用测微螺杆
沿楔方向平移。两楔的厚度之和保持不变（斜面互补），所以一片增厚多少，另一片就减薄多少：

\[
t_1+t_2=H\ (\text{常数})\quad\Longrightarrow\quad \Delta t_1=-\Delta t_2 .
\]

于是厚度\emph{差}的变化是单块厚度变化的两倍：楔平移 $x$ 时 $t_1$ 变化 $x\tan\alpha$，
$t_2$ 反向变化 $x\tan\alpha$，

\[
\Delta=\Delta n\,(t_1-t_2)\ \longrightarrow\
\Delta n\left[(t_1-t_2)+2x\tan\alpha\right].
\]

这就是课件 p72 与 p109 例题里那个因子 $2$ 的来历。

\subhdr{15.2 黑带：零延迟的等高线}

把它放在正交偏振器之间。由 §20 的公式，正交检偏器下透射强度正比于
$\sin^2(\delta/2)$，所以 $\delta=0$（即 $t_1=t_2$）的那条线\emph{全黑}，两侧对称地亮起来。
这条黑带平行于楔边；平移楔，黑带跟着移动。测微鼓轮按波长分数刻度（教学仪器 $\lambda/100$，
研究级 $\lambda/1000$），读数差就是样品的延迟。

\textbf{测量步骤}：先不放样品，记下黑带在参考线的读数；放入样品，样品的延迟把黑带推走；
再平移补偿器把黑带推回参考线；两次读数之差 = 样品的延迟（数值相等、符号相反）。

\subhdr{15.3 三条限制}

\textbf{必须单色光。}黑带的位置正比于 $\Delta/\lambda$；白光里每个波长各有一条零延迟线，
黑带糊成彩色宽带，根本读不准。

\textbf{必须加光阑。}光束有宽度，楔在光束范围内厚度不同，读数是一个平均值；光阑把采样区缩小。

\textbf{不是消色差的。}刻度严格只在标定波长成立。

\subhdr{15.4 课件例题（p72）}

\begin{center}
\small
\tabtop
\begin{tabular}{@{}p{15.2cm}@{}}
\thead{题} \\
\tabmid
巴比涅补偿器，楔角 $3^\circ$，$\Delta n=0.0091$，可动楔平移 $0.8\ \mathrm{mm}$。求 (a) 有效厚度变化；(b) 产生的延迟。\\
\tabmid
(a) $\Delta d=2x\tan\alpha=2\times0.8\ \mathrm{mm}\times\tan3^\circ=0.0838\ \mathrm{mm}$；\\
(b) $\Delta=\Delta n\,\Delta d=0.0091\times0.0838\ \mathrm{mm}=762\ \mathrm{nm}\approx1.29\lambda$（$589\ \mathrm{nm}$；课件取 $760\ \mathrm{nm}$，差在 $\tan3^\circ$ 的取位）。\\
\tabmid
\textbf{读法}：一次平移就让条纹走过一个多级次，所以测量必须相对参考位置读数，不能只看绝对刻度。\\
\end{tabular}
\tabbot
\normalsize
\end{center}

%=====================================================================
\hdr{§16}{索累--巴比涅补偿器：消色差的那一版}

两片石英楔的\emph{光轴平行}，另加一片固定厚度的石英片，其光轴与两楔垂直。两楔反向等量移动：
总厚度在变，但光束穿过的材料总量与出射方向都不变，所以视场均匀、没有条纹。

\textbf{为什么消色差}：两楔的色散与固定片的色散一阶相消，在相当宽的波段内延迟近似不随波长变。
代价是结构复杂、贵。

\begin{center}
\small
\tabtop
\begin{tabular}{@{}p{3.4cm}p{5.4cm}p{6.2cm}@{}}
\thead{} & \thead{巴比涅} & \thead{索累--巴比涅} \\
\tabmid
轴的取向 & 两楔互相垂直 & 两楔平行，另有垂直的固定片 \\
\tabmid
视场 & 有条纹（延迟随位置变） & 全场均匀 \\
\tabmid
色差 & 有，只在一个波长准 & 一阶消色差 \\
\tabmid
主要用途 & 单点或单线的精密测量 & 全场可调的「波片」 \\
\tabmid
读数方式 & 量黑带位移 & 刻度直接标定延迟 \\
\end{tabular}
\tabbot
\normalsize
\end{center}

%=====================================================================
\hdr{§17}{塞纳蒙与贝雷克：不靠楔形的两种办法}

\subhdr{17.1 塞纳蒙法：转检偏器，把延迟读成角度}

装置：样品之后放一片 $\lambda/4$ 片（快轴与起偏器透振方向平行），再放可转的检偏器。
样品把线偏振变成椭圆，$\lambda/4$ 片再把它变回线偏振，但方位转过了一个角度 $\theta$：

\[
R=\theta\cdot\frac{\lambda}{180^\circ}.
\]

\textbf{自检这个式子}：若样品是半波片（$R=\lambda/2$），代入得 $\theta=90^\circ$——半波片在 $45^\circ$
放置时正好把振动面转过 $90^\circ$，与 §10 的结论一致。

\textbf{为什么准}：细测量交给检偏器的转角（可读到 $0.1^\circ$），而不是测微螺杆；
灵敏度可达 $\lambda/1000$。代价是对波长敏感，必须配窄带滤光片（常用 $546\ \mathrm{nm}$）。
例：$\theta=30^\circ$、$\lambda=546\ \mathrm{nm}$，则 $R=30\times546/180=91\ \mathrm{nm}$。

\subhdr{17.2 贝雷克法：倾斜一片单轴晶体}

一片光轴垂直于表面的单轴晶片（常用 MgF$_2$ 或方解石），插在光路里倾斜使用。
正入射时延迟为零；一倾斜，o 光与 e 光走过的路径与折射率都变，延迟随倾角增大：

\[
\Gamma(\alpha)=d\,n_o\left[
\sqrt{1-\frac{\sin^2\alpha}{n_e^2}}-\sqrt{1-\frac{\sin^2\alpha}{n_o^2}}
\right],
\]

这就是 Nikitin--Berek 公式。倾斜角由刻度鼓轮读出，延迟由标定表换算。
它比塞纳蒙更不依赖波长，是矿物学与岩石学的常用工具。

\begin{center}
\small
\tabtop
\begin{tabular}{@{}p{3.2cm}p{5.6cm}p{6.0cm}@{}}
\thead{} & \thead{塞纳蒙} & \thead{贝雷克} \\
\tabmid
机制 & 固定 $\lambda/4$ + 转检偏器 & 倾斜单轴晶片 \\
\tabmid
调节量 & 检偏器转角 & 倾角 \\
\tabmid
波长依赖 & 强，必须单色（如 $546\ \mathrm{nm}$） & 较弱 \\
\tabmid
典型场合 & 显微镜、DIC 偏置 & 矿物学、岩石学 \\
\end{tabular}
\tabbot
\normalsize
\end{center}

%=====================================================================
\hdr{§18}{一次完整的延迟测量：五步}

\begin{enumerate}[label=\textbf{第 \arabic* 步},leftmargin=4.6em]
\item 正交放置起偏器与检偏器，确认视场全黑（这一步同时检查了系统的消光比）。
\item 放入补偿器，调到零延迟位置，记下参考读数。
\item 放入样品，把它的主轴转到与偏振器成 $45^\circ$——此时灵敏度最高（§20 的 $\sin^2 2\varphi$ 取极大）。
\item 调补偿器，把视场调回最暗（零位法），读出新读数。
\item 两读数之差即样品的延迟 $R$；再除以样品厚度得 $\Delta n=R/d$。报告时必须注明测量波长。
\end{enumerate}

\textbf{精度来自哪里}：不是来自「读一个亮场有多亮」，而是来自「把场调到多暗」。
这就是 §6.2 说过的零位法，在补偿器上体现得最彻底。

%=====================================================================
\hdr{§19}{两束正交偏振的光为什么不干涉}

\subhdr{19.1 一般公式与它的来源}

两束光振幅 $E_1$、$E_2$，相位差 $\delta$，振动方向夹角 $\theta$。叠加后的强度是

\[
I=I_1+I_2+2\sqrt{I_1I_2}\,\cos\theta\,\cos\delta .
\]

第三项就是干涉项。它的来源是两件事相乘：\emph{几何}上，两个振动只有在同一方向上的投影才能相加减，
投影因子是 $\cos\theta$；\emph{时间}上，两列波的相位差决定它们是同相加强还是反相削弱，
平均后给出 $\cos\delta$。

于是两个极端：

\begin{itemize}[leftmargin=1.8em]
\item $\theta=90^\circ$：$\cos\theta=0$，干涉项消失，$I=I_1+I_2$——\textbf{强度直接相加}，没有条纹。
\item $\theta=0^\circ$：干涉项最大，条纹对比度最好。
\end{itemize}

\subhdr{19.2 菲涅耳--阿拉戈实验（1816--1819）}

两位用实验把这件事钉死：两束偏振方向成 $90^\circ$ 的光，无论光程差怎么变，都不出现条纹；
但只要在它们后面放一个检偏器，把两束的振动投影到同一方向，条纹立刻回来。
这个实验同时证明了光是横波——纵波没有「振动方向」可言，也就不会有这种选择性干涉。

\subhdr{19.3 检偏器不是滤波器，是「让干涉成为可能」的元件}

这一点最容易被误解。波片后面的两束光确实在空间上重叠，但它们正交，叠加只是强度相加，
所以视场是均匀的亮（或均匀的暗），看不到任何图案。检偏器把两束都投影到自己透振方向上，
两个投影\emph{共线}了，才能相加减——干涉图样是检偏器「造」出来的，不是晶体里现成的。

\subhdr{19.4 干涉的四个条件}

\begin{center}
\small
\tabtop
\begin{tabular}{@{}p{3.6cm}p{11.6cm}@{}}
\thead{条件} & \thead{含义与本讲中的体现} \\
\tabmid
相干 & 来自同一光源、相位关系稳定；同一波前的两个分量自动满足 \\
\tabmid
同频率 & 波长相同；白光下每个波长各自成一套条纹 \\
\tabmid
振动不严格正交 & 必须有共同分量；靠检偏器实现 \\
\tabmid
光程差在相干长度内 & 钠灯几毫米，激光几米，白光只有几微米——所以白光下只能用薄片 \\
\end{tabular}
\tabbot
\normalsize
\end{center}

%=====================================================================
\hdr{§20}{核心公式：正交检偏器之间的波片}

\subhdr{20.1 装置与推导}

起偏器透振方向取 $x$；波片快轴与 $x$ 成 $\varphi$；检偏器与起偏器正交（沿 $y$）。
入射振幅 $E_0$，过起偏器后 $\vect E=(E_0,0)$。

第一步，分解到波片的两条轴上。入射振动沿 $x$，而 $x$ 相对快轴的方位是 $-\varphi$（不是 $+\varphi$），
所以沿快轴、慢轴的系数分别是 $E_0\cos\varphi$ 与 $-E_0\sin\varphi$。
第二步，过片：慢轴分量落后 $\delta$，系数再乘 $e^{-i\delta}$。
第三步，把两个分量投影到检偏器方向 $y$。快轴单位矢量 $(\cos\varphi,\sin\varphi)$ 与慢轴单位矢量
$(-\sin\varphi,\cos\varphi)$ 在 $y$ 上的投影分别是 $\sin\varphi$ 与 $\cos\varphi$，于是

\[
E_y=E_0\cos\varphi\cdot\sin\varphi+\left(-E_0\sin\varphi\,e^{-i\delta}\right)\cdot\cos\varphi
=E_0\sin\varphi\cos\varphi\left(1-e^{-i\delta}\right).
\]

\textbf{用一个特例检查符号}：$\varphi=45^\circ$、半波片 $\delta=\pi$，上式给出
$E_y=E_0\cdot\tfrac12\cdot 2=E_0$，即交叉检偏器全亮——半波片把振动面转了 $90^\circ$，正该如此。
若符号反了（写成 $e^{-i\delta}-1$），这一步会给出零，与实验矛盾。

取模平方：

\[
I=|E_y|^2=E_0^2\sin^2\varphi\cos^2\varphi\left|1-e^{-i\delta}\right|^2
=E_0^2\frac{\sin^2 2\varphi}{4}\cdot 4\sin^2\frac{\delta}{2},
\]

其中用了 $|1-e^{-i\delta}|^2=2-2\cos\delta=4\sin^2(\delta/2)$。于是

\[
\boxed{\;I=I_0\sin^2 2\varphi\;\sin^2\frac{\delta}{2}\;}
\qquad(\text{正交检偏器}),\qquad \delta=\frac{2\pi}{\lambda}\Delta n\,d .
\]

\subhdr{20.2 平行检偏器}

若检偏器与起偏器平行，投影取 $x$ 方向，同样的推导给出

\[
I=I_0\left[1-\sin^2 2\varphi\;\sin^2\frac{\delta}{2}\right].
\]

两式互补：同一块片子，正交时暗的地方，平行时就亮。这一点在实验里很好用——
把检偏器转 $90^\circ$，图案应当明暗互换，否则说明装置里有杂散光。

\subhdr{20.3 两个因子各管什么}

\textbf{$\sin^2 2\varphi$ 是几何因子}：它只问「两个分量各有多少被投影到检偏器方向上」。
$\varphi=0^\circ$ 或 $90^\circ$ 时为零（波片的轴与偏振器对齐，等于没放片子）；
$\varphi=45^\circ$ 时为 $1$（最大对比度）。

\textbf{$\sin^2(\delta/2)$ 是物理因子}：它只问「波片制造了多大的相位差」。
$\delta=0$ 或 $2\pi$ 时为零（全波片、或者根本没放片子）；
$\delta=\pi$ 时为 $1$（半波片把振动面转到与检偏器平行，全亮）。

\subhdr{20.4 三个必会的工作例（课件 p100）}

\textbf{半波片 $45^\circ$ 放置。}$\delta=\pi$，$\sin^2(\delta/2)=1$；$\varphi=45^\circ$，$\sin^2 2\varphi=1$。
于是 $I=I_0$：正交偏振器之间全亮。原因不是「波片发光」，而是半波片把振动面转了 $90^\circ$，
正好对上检偏器。同一片放在平行偏振器之间则是 $I=0$。

\textbf{四分之一波片 $45^\circ$ 放置。}$\delta=\pi/2$，$\sin^2(\pi/4)=1/2$；$\varphi=45^\circ$。
于是 $I=I_0/2$，而且\emph{转动检偏器强度不变}——出射是圆偏振，任何方向的投影都是一半。
这是判断「是否真的得到圆偏振」的标准检验。

\textbf{全波片 $45^\circ$ 放置。}$\delta=2\pi$，$\sin^2\pi=0$，$I=0$：偏振态被完全复原，
检偏器把它挡掉了。只有在白光下它才「显形」，因为其他波长不满足 $\delta=2\pi$。

%=====================================================================
\hdr{§21}{色偏振：延迟就是一把彩色的尺子}

\subhdr{21.1 消光波长}

白光入射时，$\sin^2(\delta/2)=0$ 的条件是 $\delta=2\pi m$，即

\[
\Delta=m\lambda\quad\Longrightarrow\quad \lambda_m=\frac{\Delta}{m},\qquad m=1,2,3,\dots
\]

也就是说，一块延迟为 $\Delta$ 的片子会把 $\Delta,\ \Delta/2,\ \Delta/3,\dots$ 这些波长从透射谱里挖掉。
挖掉的是一组窄带，剩下的是「被剪过的白光」——看上去有颜色。这就是色偏振，颜色来自干涉，不来自吸收。

\subhdr{21.2 牛顿色序}

\begin{center}
\small
\tabtop
\begin{tabular}{@{}p{2.2cm}p{4.2cm}p{8.8cm}@{}}
\thead{级次} & \thead{延迟（nm）} & \thead{典型颜色} \\
\tabmid
一级 & $0$--$550$ & 灰 $\to$ 白 $\to$ 黄 $\to$ 红 \\
\tabmid
二级 & $550$--$1100$ & 蓝 $\to$ 绿 $\to$ 黄 $\to$ 红 \\
\tabmid
三级 & $1100$--$1650$ & 蓝 $\to$ 绿 $\to$ 黄 $\to$ 红 \\
\tabmid
更高级 & $>1650$ & 淡而发白的粉彩色 \\
\end{tabular}
\tabbot
\normalsize
\end{center}

\subhdr{21.3 课件 p95 的一处口径需要核对}

课件举了 $\Delta=30\ \mu\mathrm{m}\times0.0091=273\ \mathrm{nm}$ 的例子，说可见区「只有一处消光，
在 $546\ \mathrm{nm}$ 附近」。按它自己给的 $\lambda_m=\Delta/m$ 核对：$273/1=273\ \mathrm{nm}$（紫外），
$273/2=136.5\ \mathrm{nm}$（更远），\emph{可见区里一处消光也没有}。
$546\ \mathrm{nm}$ 处 $\Delta/\lambda=1/2$，代进公式得 $\sin^2(\pi/2)=1$——那是\emph{极大}，不是消光。

所以「颜色偏淡黄绿」的结论是对的，理由要改：$546\ \mathrm{nm}$ 附近透过率最大，
红光（$700\ \mathrm{nm}$）约 $0.89$、蓝光（$400\ \mathrm{nm}$）约 $0.71$（都归一到 $546\ \mathrm{nm}$），
加权之后偏黄绿。写作业时按「极大在 $546\ \mathrm{nm}$」说，不要跟着写成「消光在 $546\ \mathrm{nm}$」。

\subhdr{21.4 为什么厚片看起来发白}

延迟越大，可见区里的消光带越多、越密。等到几十条带挤在 $400$--$700\ \mathrm{nm}$ 里，
人眼已经把光谱平均掉了，看到的就是「高级白」。所以观察色偏振要用薄片：
矿物薄片的标准厚度是 $30\ \mu\mathrm{m}$，道理就在这里。

%=====================================================================
\hdr{§22}{干涉图：等色线与等旋线}

\subhdr{22.1 平行光与会聚光}

\textbf{平行光（orthoscopic）}：所有光线以同一角度穿过晶体，延迟处处相同，视场是均匀的一种颜色；
把片子倾斜，颜色跟着变（有效延迟变了）。若片子是楔形，出现平行于楔边的直条纹，
每条条纹是一条等延迟线（等色线）。

\textbf{会聚光（conoscopic）}：光线以各种角度穿过晶体，每一点的延迟都不同，于是视场里出现
一族彩色曲线（等色线，isochromats）加一组暗刷（等旋线，isogyres）。

\subhdr{22.2 十字与它的旋转}

等旋线出现在「晶体内的振动方向与偏振器方向重合」的地方，所以单轴晶体在会聚光下出现一个
以光轴为中心的\emph{黑十字}。转动载物台：

\begin{itemize}[leftmargin=1.8em]
\item 单轴晶体：十字保持居中，不分裂——这是单轴的判据。
\item 双轴晶体：十字分裂成两条双曲线，并随台转动而旋转——这是双轴的判据。
\end{itemize}

\subhdr{22.3 怎么判双折射的正负}

在光路里插入一片已知延迟的 $\lambda/4$ 片或石膏板（一级红），观察光轴附近颜色的升降：
延迟相加则颜色升高（蓝变高），相减则降低。这一个动作就把「慢轴是 o 还是 e」定下来了，
也就是正晶体还是负晶体。

%=====================================================================
\hdr{§23}{三个把干涉用起来的地方}

\subhdr{23.1 Lyot 滤光片：用一叠波片选出极窄的一条谱线}

把一串波片放在平行偏振器之间，厚度依次加倍（$d,2d,4d,\dots$）。每一级都把「不该通过的波长」
再筛一遍，结果是：

\begin{itemize}[leftmargin=1.8em]
\item 最薄的那片决定\emph{自由光谱范围}（能隔多远出现下一个通带）；
\item 整叠片子的总数决定\emph{通带宽度}：六片可以做到 $0.1\ \mathrm{nm}$ 以下。
\end{itemize}

太阳望远镜用它隔离 H$\alpha$ 线（$656.3\ \mathrm{nm}$），卫星遥感用它做高光谱成像。

\subhdr{23.2 光弹应力分析：把应力场变成条纹}

透明模型受力后产生应力双折射，延迟与主应力差成正比：

\[
\Delta=C\,d\,(\sigma_1-\sigma_2)\quad\Longrightarrow\quad
\sigma_1-\sigma_2=\frac{n\lambda}{d\,C},
\]

其中 $C$ 是应力光学系数，$n=\Delta/\lambda$ 是条纹级数。把模型放进正交偏振器之间，
每一条条纹就是「主应力差相等」的等高线，孔的边缘条纹越密，说明应力集中越严重。

\textbf{为什么要用圆偏振光场}：普通正交偏振器下还会出现等旋线（与主应力\emph{方向}有关），
和等色线混在一起看不清。在样品前后各加一片 $\lambda/4$ 片（圆偏振器），等旋线被去掉，
只剩等色线，条纹级数可以直接数。

\subhdr{23.3 电光与磁光调制器}

\textbf{Pockels 盒}：加电压在晶体里诱导出正比于电场的双折射，$\Delta n\propto E$，晶体变成一个
延迟可电调的波片；夹在正交偏振器之间就成了光开关，响应可到纳秒。

\textbf{Kerr 盒}：$\Delta n\propto E^2$，需要更高电压，但响应快、口径大。

\textbf{Faraday 旋转器}：磁场让振动面旋转，$\theta=VBL$；配上偏振器就是光隔离器——
激光打回头的光被挡住，保护光源。

\textbf{共同点}：都是「用电或磁造一个可控的双折射/旋转」，再由检偏器把状态变成强度。
这正是 §19 的逻辑。

%=====================================================================
\hdr{§24}{课件里的三道例题（教学示范）}

这三题课件已经给出答案，所以它们在这里是\emph{示范}，不是留给你做的练习：看的是每一步「为什么这样代」。

\subhdr{24.1 例题一：尼科尔棱镜 + 马吕斯定律（p107）}

\qstem{（课件 p107）自然光 $I_0=100\ \mathrm{W/m^2}$ 射到尼科尔棱镜上，透射光再过一个检偏棱镜，
两者主截面成 $30^\circ$。求 (a) 第一块之后；(b) 检偏器之后；(c) 检偏器转到 $90^\circ$ 时。}

\textbf{(a)} 偏振棱镜对自然光只放行一个分量：

\[
I_1=\frac{I_0}{2}=\frac{100}{2}=50\ \mathrm{W/m^2}.
\]

\textbf{(b)} 入射已是线偏振，用 Malus 定律，夹角取\emph{两块棱镜之间}的 $30^\circ$：

\[
I_2=I_1\cos^2 30^\circ=50\times0.75=37.5\ \mathrm{W/m^2}.
\]

\textbf{(c)} 正交时 $I_3=I_1\cos^2 90^\circ=0$。实际测到的不是绝对零，而是棱镜的漏光——
这个残余强度与 $I_1$ 之比就是消光比。

\textbf{读法}：整条链上只有第一次是「减半」，后面每一步都是 $\cos^2$。把顺序搞反（在 $100$ 上直接乘
$\cos^2 30^\circ$）是最常见的失分点。

\subhdr{24.2 例题二：石英四分之一波片（p108）}

\qstem{（课件 p108）石英片平行光轴切割，厚 $d=16.2\ \mu\mathrm{m}$；$\lambda=589\ \mathrm{nm}$，
$n_e-n_o=+0.0091$（石英是正晶体）。(a) 求延迟（纳米与度）；(b) 判断是什么片；(c) $45^\circ$
线偏振入射时出射是什么态？}

\textbf{(a)} 光程差与相位差：

\[
\Delta=\Delta n\,d=0.0091\times16.2\ \mu\mathrm{m}=0.1474\ \mu\mathrm{m}\approx147\ \mathrm{nm},
\]
\[
\delta=360^\circ\times\frac{\Delta}{\lambda}=360^\circ\times\frac{147}{589}\approx89.9^\circ\approx90^\circ .
\]

\textbf{(b)} 与 $\lambda/4=589/4=147.25\ \mathrm{nm}$ 相比几乎相等，所以是四分之一波片；
厚度取的是最小可能值，所以是\textbf{零级}片。

\textbf{(c)} $45^\circ$ 入射 → 两分量等幅；延迟 $90^\circ$ → 出射\textbf{圆偏振}。
严格说是「略偏的椭圆」，因为 $89.9^\circ\ne90^\circ$——课件特意点出这一点，说明它在意的是
「设计值」与「实际值」的差别。

\subhdr{24.3 例题三：巴比涅补偿器测双折射（p109）}

\qstem{（课件 p109）巴比涅补偿器置于正交偏振器之间，$\lambda=589\ \mathrm{nm}$，石英楔角 $2.5^\circ$，
$\Delta n=0.0091$。放入厚 $2.0\ \mathrm{mm}$ 的受力塑料样品后，把楔平移 $1.0\ \mathrm{mm}$ 使黑带复位。
求 (a) 补偿器引入的延迟；(b) 样品的双折射。}

\textbf{(a)} 两楔厚度一增一减，有效厚度变化是单块的两倍：

\[
\Delta=\Delta n\,(2x\tan\alpha)
=0.0091\times2\times1.0\ \mathrm{mm}\times\tan2.5^\circ
=0.0091\times0.0873\ \mathrm{mm}\approx795\ \mathrm{nm},
\]

约 $1.35$ 个波长（$589\ \mathrm{nm}$）。

\textbf{(b)} 补偿器与样品等值反向，所以样品的延迟也是 $795\ \mathrm{nm}$；除以样品厚度：

\[
\Delta n_{\text{样品}}=\frac{\Delta}{d}=\frac{795\ \mathrm{nm}}{2.0\ \mathrm{mm}}\approx4.0\times10^{-4}.
\]

\textbf{读法}：$4\times10^{-4}$ 是「中等应力塑料」的典型值。补偿器的意义就在这里——
把一个看不见的相位差变成一个带单位的数。

%=====================================================================
\hdr{§25}{作业}

\subhdr{25.1 本讲作业（课件 p110）}

\practice{P4}{课件 p110 Problem 1}{自然光强 $I_0$ 依次通过两块偏振棱镜，两透振方向夹角 $45^\circ$。
(a) 有几分之几的 $I_0$ 从第二块出来？(b) 若第一块是尼科尔棱镜、第二块转到 $60^\circ$，出射又是多少？
(c) 若在两块之间插入第三块、其透振方向与第一块成 $30^\circ$，结果如何？解释。}
{自然光减半；之后每片都用 Malus，角度取\emph{相邻}两片；三片串联不是强度的线性叠加}{ans:P4}

\practice{P5}{课件 p110 Problem 2}{$633\ \mathrm{nm}$ 的半波片用石英制作（$\Delta n=0.0091$）。
(a) 求最小厚度（零级）；(b) 振动方向与快轴成 $20^\circ$ 的线偏振光通过后，出射方位是多少？
(c) 圆偏振光入射，出射是什么态？(d) 解释为什么这片波片在 $450\ \mathrm{nm}$ 处不能当半波片用。}
{$\lambda/2$ 厚度 $\lambda/(2\Delta n)$；方位公式 $\alpha_{\text{out}}=2\theta-\alpha_{\text{in}}$；
延迟与波长成反比}{ans:P5}

\practice{P6}{课件 p110 Problem 3}{$d=30\ \mu\mathrm{m}$、$\Delta n=0.0091$ 的双折射片放在正交偏振器之间，
快轴与起偏器成 $45^\circ$，用白光照明。(a) 求延迟；(b) 可见区 $400$--$700\ \mathrm{nm}$ 内哪些波长被消光？
(c) 说明透射光的颜色，并解释为什么换用其中某个波长的单色光时视场是黑的；(d) 若片子略有楔形，视场是什么样？}
{$\Delta=\Delta n\,d$；消光条件 $\Delta=m\lambda$；$I=I_0\sin^2 2\varphi\sin^2(\pi\Delta/\lambda)$}{ans:P6}

\subhdr{25.2 上一讲作业（课件 p60，上次已交）}

\practice{P1}{课件 p60 Problem 1}{自然光 $I_0=80\ \mathrm{W/m^2}$ 过第一片线偏振片，再过一个检偏器，
其透振方向与第一片成 $\theta=60^\circ$。求 (a) 第一片后；(b) 第二片后；(c) 第二片后为 $10\ \mathrm{W/m^2}$
时 $\theta$ 是多少。}{自然光过第一片减半；之后用 Malus；第三问反解 $\cos^2\theta$，注意两个解}{ans:P1}

\practice{P2}{课件 p60 Problem 2}{光在空气中（$n_1=1.00$）射到玻璃表面（$n_2=1.50$）。
(a) 求布儒斯特角；(b) 自然光入射时，$\theta_B$ 处反射光的偏振态；(c) 若入射是 p 偏振，$\theta_B$ 处反射强度是多少。}
{$\tan\theta_B=n_2/n_1$；$\theta_B$ 处 $R_p=0$、$\theta_B+\theta_t=90^\circ$}{ans:P2}

\practice{P3}{课件 p60 Problem 3}{双折射晶体 $n_o=1.658$、$n_e=1.486$，$\lambda=500\ \mathrm{nm}$，
光垂直光轴传播。(a) 求四分之一波片的最小厚度；(b) 厚 $10\ \mu\mathrm{m}$ 的片子产生多少度相位差；
(c) 与光轴成 $45^\circ$ 的线偏振光过 (a) 的波片后是什么偏振态。}
{$\Delta\varphi=2\pi|n_o-n_e|d/\lambda$；$\lambda/4$ 厚度；$45^\circ$ 等幅 + $90^\circ$ 相差 = 圆偏振}{ans:P3}

\clearpage
\appendix

%=====================================================================
\hdr{附录 A}{公式速查}

\begin{center}
\small
\tabtop
\begin{tabular}{@{}p{6.4cm}p{8.6cm}@{}}
\thead{公式} & \thead{条件与说明} \\
\tabmid
$I=\dfrac{I_0}{2}$ & 自然光过第一片偏振器（棱镜或偏振片）\\
\tabmid
$I=I_0\cos^2\alpha$ & 线偏振入射；$\alpha$ 是相邻两片的夹角 \\
\tabmid
$\tan\theta_B=\dfrac{n_2}{n_1}$ & 电介质界面；此时 $\theta_B+\theta_t=90^\circ$，$R_p=0$ \\
\tabmid
$\sin\theta_c=\dfrac{n_b}{n_o}$ & 尼科尔/格兰棱镜里 o 光在胶层上的临界角 \\
\tabmid
$\Delta n=|n_e-n_o|$ & 方解石 $0.172$，石英 $0.009$ \\
\tabmid
$\Delta=\Delta n\,d$，$\delta=\dfrac{2\pi}{\lambda}\Delta$ & 波片的光程差与相位差 \\
\tabmid
$d_{\lambda/4}=\dfrac{\lambda}{4\Delta n}$，$d_{\lambda/2}=\dfrac{\lambda}{2\Delta n}$ & 零级波片厚度 \\
\tabmid
$\alpha_{\text{out}}=2\theta-\alpha_{\text{in}}$ & 半波片：出射方位关于快轴镜像 \\
\tabmid
$I=I_0\sin^2 2\varphi\,\sin^2\dfrac{\delta}{2}$ & 正交检偏器之间的波片 \\
\tabmid
$I=I_0\left[1-\sin^2 2\varphi\,\sin^2\dfrac{\delta}{2}\right]$ & 平行检偏器之间 \\
\tabmid
$\lambda_m=\dfrac{\Delta}{m}$ & 色偏振的消光波长 \\
\tabmid
$R=\theta\cdot\dfrac{\lambda}{180^\circ}$ & 塞纳蒙法，$\theta$ 为检偏器转角（度）\\
\tabmid
$\sigma_1-\sigma_2=\dfrac{n\lambda}{d\,C}$ & 光弹：条纹级数 $n$ 换主应力差 \\
\end{tabular}
\tabbot
\normalsize
\end{center}

%=====================================================================
\hdr{附录 B}{本讲作业（p110）详解}

\answerstart{P4}{课件 p110 Problem 1}{自然光 $I_0$ 过两块透振方向成 $45^\circ$ 的偏振棱镜；
(a) 出射占几分之几；(b) 第二块转到 $60^\circ$；(c) 中间插入一块与第一块成 $30^\circ$ 的棱镜。}

\textbf{(a)} 自然光过第一块减半，第二块用 Malus：

\[
I_2=\frac{I_0}{2}\cos^2 45^\circ=\frac{I_0}{2}\times\frac12=\frac{I_0}{4},
\]

即 $\dfrac14$（$25\%$）。

\textbf{(b)} 只换角度，方法不变：

\[
I_2'=\frac{I_0}{2}\cos^2 60^\circ=\frac{I_0}{2}\times\frac14=\frac{I_0}{8}=12.5\%.
\]

\textbf{(c)} 三块串联，角度要一段一段取：第一块到第三块是 $30^\circ$，第三块到第二块是 $45^\circ-30^\circ=15^\circ$：

\[
I_3=\frac{I_0}{2}\cos^2 30^\circ\cos^2 15^\circ
=\frac{I_0}{2}\times0.75\times0.9330\approx0.350\,I_0 .
\]

\textbf{为什么插一片反而更亮}：两块正交时本来全黑（$\cos^2 90^\circ=0$）；中间那片把振动方向
「转」到 $30^\circ$，于是第一片到中间片有 $30^\circ$ 的投影，中间片到第二片有 $15^\circ$ 的投影，
两次都不为零，乘积就不是零。\textbf{强度在偏振片链上不能相加}，这是 Malus 定律最容易出反直觉结论的地方。

\textbf{自检}：$0.350<0.25$ 吗？不——$0.350>0.25$，插片后比两块成 $45^\circ$ 时还亮一点。
这并不矛盾：两块 $45^\circ$ 时是 $\cos^2 45^\circ=0.5$，插片后是 $0.75\times0.933=0.70$，确实更大。

%---------------------------------------------------------------------
\answerstart{P5}{课件 p110 Problem 2}{石英半波片，$\lambda=633\ \mathrm{nm}$，$\Delta n=0.0091$。
(a) 最小厚度；(b) 方位 $20^\circ$ 的线偏振入射；(c) 圆偏振入射；(d) $450\ \mathrm{nm}$ 处为何失效。}

\textbf{(a)} 半波条件 $\Delta n\,d=\lambda/2$：

\[
d=\frac{\lambda}{2\Delta n}=\frac{633\ \mathrm{nm}}{2\times0.0091}
=\frac{633}{0.0182}\ \mathrm{nm}\approx34\,780\ \mathrm{nm}\approx34.8\ \mu\mathrm{m}.
\]

取 $m=0$ 就是最薄的零级片。

\textbf{(b)} 记快轴为 $0^\circ$，入射方位 $\alpha_{\text{in}}=20^\circ$：

\[
\alpha_{\text{out}}=2\times0^\circ-20^\circ=-20^\circ .
\]

即出射仍是线偏振，方向关于快轴对称——从 $+20^\circ$ 翻到 $-20^\circ$。
若要把方位转到别的角度，把波片转到 $\theta=(\alpha_{\text{in}}+\alpha_{\text{out}})/2$ 即可。

\textbf{(c)} 圆偏振过 $\lambda/2$ 片\textbf{换手}：左旋变右旋，右旋变左旋。振幅不变、椭圆率不变。

\textbf{(d)} $450\ \mathrm{nm}$ 处的延迟不再是半个波长。用同一片算：

\[
\frac{\Delta}{\lambda}=\frac{0.0091\times34\,780\ \mathrm{nm}}{450\ \mathrm{nm}}
=\frac{316.5}{450}\approx0.703,
\]

即 $\delta\approx253^\circ$，离 $180^\circ$ 差了 $73^\circ$。出射成为椭圆偏振而不是线偏振，
所以它在这条谱线上不是半波片。这也说明\textbf{波片必须按工作波长选}。

%---------------------------------------------------------------------
\answerstart{P6}{课件 p110 Problem 3}{$d=30\ \mu\mathrm{m}$、$\Delta n=0.0091$ 的片子，正交偏振器之间，
快轴成 $45^\circ$，白光照明。}

\textbf{(a)}

\[
\Delta=\Delta n\,d=0.0091\times30\ \mu\mathrm{m}=0.273\ \mu\mathrm{m}=273\ \mathrm{nm}.
\]

\textbf{(b)} 正交偏振器下消光条件 $\Delta=m\lambda$，即 $\lambda_m=\Delta/m$：

\[
\lambda=273\ \mathrm{nm}\ (m=1),\quad 136.5\ \mathrm{nm}\ (m=2),\quad 91\ \mathrm{nm}\ (m=3),\dots
\]

这些全在紫外，\textbf{可见区 $400$--$700\ \mathrm{nm}$ 内没有消光波长}。
顺便把 $546\ \mathrm{nm}$ 代进去：$\Delta/\lambda=0.5$，$\sin^2(\pi/2)=1$——那是极大值。

\textbf{(c)} 透射光偏\textbf{淡黄绿}：$546\ \mathrm{nm}$ 附近透过率最大（为 $1$），
$700\ \mathrm{nm}$ 处约 $0.89$、$400\ \mathrm{nm}$ 处约 $0.71$，红蓝都被削掉一部分，剩下偏绿。
换用 $\lambda=273\ \mathrm{nm}$ 的单色光时，$\Delta=\lambda$，$\sin^2(\delta/2)=\sin^2\pi=0$，视场全黑——
这就是「单色光下能看见消光」而白光下只能看见颜色的区别。

\textbf{(d)} 楔形意味着厚度沿一个方向线性变化，于是 $\Delta$ 也线性变化，条纹是\emph{平行于楔边的一系列
直条纹}（等色线）；每条条纹对应一个固定的延迟（也就对应一个固定的颜色）。若用单色光，
看到的是明暗相间的直条纹；用白光，看到的是彩色条纹带，颜色按牛顿色序排列。

\textbf{自检}：$273\ \mathrm{nm}$ 恰在 $546\ \mathrm{nm}$ 的一半处，所以 $546\ \mathrm{nm}$ 是「半个波长延迟」
对应的极大——这也解释了为什么这片在白光下颜色很淡：整个可见区都离消光条件很远。

%=====================================================================
\hdr{附录 C}{上一讲作业三题（课件 p60）详解}

\answerstart{P1}{课件 p60 Problem 1}{自然光 $I_0=80\ \mathrm{W/m^2}$；两片透振方向夹角 $\theta=60^\circ$。
求第一片后、第二片后，以及第二片后为 $10\ \mathrm{W/m^2}$ 时的角度。}

\textbf{已知}：$I_0=80\ \mathrm{W/m^2}$（自然光），$\theta=60^\circ$。
\textbf{所用公式}：$I_1=I_0/2$；$I_2=I_1\cos^2\theta$。

\textbf{(a)} 第一片：自然光两个正交分量各占一半，偏振片只放行一个：

\[
I_1=\frac{80}{2}=40\ \mathrm{W/m^2}.
\]

\textbf{(b)} 第二片：入射已是线偏振，夹角 $60^\circ$：

\[
I_2=I_1\cos^2 60^\circ=40\times\left(\frac12\right)^2=10\ \mathrm{W/m^2}.
\]

\textbf{(c)} 反解。第一片后的 $40\ \mathrm{W/m^2}$ 与角度无关，是固定的：

\[
40\cos^2\theta=10\ \Longrightarrow\ \cos^2\theta=\frac14
\ \Longrightarrow\ \cos\theta=\pm\frac12
\ \Longrightarrow\ \theta=60^\circ\ \text{或}\ 120^\circ .
\]

两个解对应检偏器轴关于第一片轴对称的两个取向；再加 $180^\circ$ 回到同一条轴，
所以 $0^\circ\le\theta<180^\circ$ 内就这两个。

\textbf{答案}：(a) $40\ \mathrm{W/m^2}$；(b) $10\ \mathrm{W/m^2}$；(c) $60^\circ$ 或 $120^\circ$。

\textbf{自检}：(c) 的解里含 (b) 用过的 $60^\circ$——(c) 本来就是 (b) 的反问题，数值自洽。

\textbf{易错}：在 $80$ 上直接乘 $\cos^2 60^\circ$（得 $20$）——那是把自然光当成了线偏振，
第一片减半这一步没做。

%---------------------------------------------------------------------
\answerstart{P2}{课件 p60 Problem 2}{空气到玻璃：$n_1=1.00$，$n_2=1.50$。
求布儒斯特角、自然光入射时 $\theta_B$ 处反射光的偏振态、以及 p 偏振入射时的反射强度。}

\textbf{已知}：$n_1=1.00$，$n_2=1.50$。
\textbf{所用公式}：$\tan\theta_B=n_2/n_1$；$\theta_B+\theta_t=90^\circ$；$R_p=|r_p|^2$。

\textbf{(a)}

\[
\tan\theta_B=\frac{1.50}{1.00}=1.50
\quad\Longrightarrow\quad
\theta_B=\arctan1.50=56.31^\circ\approx56.3^\circ .
\]

\textbf{(b)} 在 $\theta_B$ 上 $r_p=0$，反射光里只剩 s 分量：\textbf{完全线偏振，振动垂直于入射面}。
注意入射光本身是非偏振的——「非偏振入射、完全偏振反射」正是布儒斯特角最反直觉的地方。

\textbf{(c)} p 偏振入射时 $r_p=0$，故

\[
I_{\text{反射}}=I_0\,R_p=I_0|r_p|^2=0 .
\]

这是精确的零，不是「反射很弱」。

\textbf{答案}：(a) $56.3^\circ$；(b) $100\%$ 线偏振，垂直于入射面；(c) $0$。

\textbf{自检}（Snell）：$\sin56.31^\circ=0.8321$，$1.50\sin33.69^\circ=1.50\times0.5547=0.8321$，
两边相等，说明 $\theta_B+\theta_t=90^\circ$ 成立。

\textbf{为什么是 $\tan$ 不是 $\sin$}：把 $\theta_t=90^\circ-\theta_B$ 代入 $n_1\sin\theta_B=n_2\sin\theta_t$，
右端变成 $n_2\cos\theta_B$，两边相除即得 $\tan\theta_B=n_2/n_1$。

%---------------------------------------------------------------------
\answerstart{P3}{课件 p60 Problem 3}{$n_o=1.658$、$n_e=1.486$、$\lambda=500\ \mathrm{nm}$，光垂直光轴传播。
求四分之一波片的最小厚度、$10\ \mu\mathrm{m}$ 厚片的相位差（度）、以及 $45^\circ$ 线偏振过该波片后的偏振态。}

\textbf{已知}：$n_o=1.658$，$n_e=1.486$，$\lambda=500\ \mathrm{nm}$，$|n_o-n_e|=0.172$。
\textbf{所用公式}：$\Delta\varphi=\dfrac{2\pi}{\lambda}|n_o-n_e|d$；$d_{\lambda/4}=\dfrac{\lambda}{4|n_o-n_e|}$。

\textbf{(a)}

\[
d_{\min}=\frac{500\ \mathrm{nm}}{4\times0.172}=\frac{500}{0.688}\ \mathrm{nm}
=726.7\ \mathrm{nm}\approx0.727\ \mu\mathrm{m}.
\]

\textbf{(b)} $d=10\ \mu\mathrm{m}=1.0\times10^4\ \mathrm{nm}$：

\[
\Delta\varphi=\frac{2\pi}{500}\times0.172\times1.0\times10^4=2\pi\times3.44=21.61\ \mathrm{rad},
\]
\[
\Delta\varphi=3.44\times360^\circ=1238.4^\circ\equiv158.4^\circ\ (\mathrm{mod}\ 360^\circ).
\]

\textbf{(c)} 与光轴成 $45^\circ$ 的线偏振，两个分量\textbf{等幅}；四分之一波片再给 $90^\circ$ 相位差，
所以出射是\textbf{圆偏振}。

\textbf{答案}：(a) $726.7\ \mathrm{nm}\approx0.727\ \mu\mathrm{m}$；
(b) $1238.4^\circ$（等价于 $158.4^\circ$；$21.61\ \mathrm{rad}$，即 $3.44$ 个波长）；
(c) 圆偏振（等幅、相差 $90^\circ$）。

\textbf{自检}：$3.44$ 个波长的延迟介于 $3$ 与 $4$ 之间，说明它不是 $\lambda/4$ 也不是 $\lambda/2$ 片；
而 (a) 算出的厚度正好对应 $0.25$ 个波长，与「$45^\circ$ 入射得到圆偏振」的条件吻合。

\textbf{关于旋向}：若老师要求写左右旋，用琼斯矢量说明约定。取 e 轴为 $x$、o 轴为 $y$、
时间因子 $e^{i(kz-\omega t)}$，则

\[
\vect J_{\text{out}}=\frac{1}{\sqrt2}\begin{pmatrix}1\\ e^{i\pi/2}\end{pmatrix}
=\frac{1}{\sqrt2}\begin{pmatrix}1\\ i\end{pmatrix},
\]

按本课程讲义 §1.4 的约定记为右旋；不同教材的约定可能相反，所以答题时把约定写在旁边，
核心结论「等幅 + $90^\circ$」与约定无关。

\end{document}







